\chapter{Discussion and Future Directions}\label{chap:discussion}
% ---- TAKEAWAYS ----
% - Handle label noise realistically: ask where it comes from and what shape it has, not only how much there is.
% - Real assay noise is not only Gaussian, and Gaussian noise is the hardest case to correct.
% - No single method fits every dataset, for building the hybrid or for learning with noisy labels.
% - Simple is often better: tree ensembles under noise, fewer features, and neighbours on a small task-relevant feature set.
% - Correcting labels without ground truth is hard; prefer LNL methods that change training over noise mitigation, except for noise of known structure.
% - Next steps: classification and flipped labels, hybrids as inputs to generative models, and a faster search.
% ---- END TAKEAWAYS ----
% ---- NOTES: Discussion and Future Directions ----
% Takeaway: before choosing a method, ask where the noise comes from and what shape it has.
% Paragraph 1: handle noise realistically (short opening, framing only, no new results).
% - Real assay labels carry shared offsets per laboratory \citep{Bentz2013} and labels clipped at detection limits \citep{Svensson2025} (see "Noise in molecular data", Chapter~\ref{chap:background}).
% - Gaussian noise is the hardest case: every label is off by an unknown amount in an unknown direction (argued in "By noise type", Chapter~\ref{chap:similarity}).
% - Cleaning labels can leave a model worse than the uncorrected labels would (evidence in "By noise type").
% - What to ask of a dataset (pooled from many laboratories, one controlled experiment, unknown source): one sentence, referring to the practical takeaways in "By noise type".
% Paragraph 2: what the chapter does (one or two sentences).
% - A summary of each chapter, then two arguments: no single method fits every dataset, and simple is often better.
% - Then the limitations and the future directions.
% ---- END NOTES ----



\section{Summary}
% ---- NOTES: Summary ----
% Takeaway: each chapter's question and its answer, one paragraph per chapter, in the order the questions were asked in Section~\ref{sec:intro:questions}.
% Paragraph 1: the problem and the tools (Chapters~\ref{chap:intro} and~\ref{chap:background}).
% - QSAR models are trained on labels that carry experimental noise, and most work assumes that noise is Gaussian.
% - Chapter~\ref{chap:background} defined the representations, models, noise types, NoiseInject and the evaluation. One sentence.
% Paragraph 2: noise robustness of representations and models (Chapter~\ref{chap:robustness}).
% - Asked: the four questions on representation, model, noise type and uncertainty under label noise.
% - Accuracy: model and representation explain comparable shares of the variance in R2 (QM9 HOMO-LUMO gap, Gaussian noise level 1.0: model 29.1%, representation 26.6%).
% - Robustness (\aucnorm): the model dominates. Under Gaussian noise the model explains 49.5% of its variance and the representation 9.2%.
% - Noise type: changing the shape of the noise barely changed performance. Grouped-shifted and censored noise cost the most.
% - Median \aucnorm was 0.813 under censoring against 0.951 under Gaussian noise, over the same five pairings of model and representation. Censoring was run on those five only.
% - Uncertainty: mean predicted uncertainty rose with noise, but no model's per-molecule uncertainty found the corrupted labels.
% - Under six of the seven noise conditions, the Spearman correlation between a molecule's predicted uncertainty and its injected noise stayed between -0.16 and 0.16.
% - Under censoring the models were most confident on the clipped labels.
% Paragraph 3: why similarity fails on standard representations (Chapter~\ref{chap:flaws}).
% - Asked: do standard representations, and distances computed over them, fail on a given dataset?
% - Many features of a general-purpose fingerprint do not vary, or do not relate to the property, on a given dataset. A distance sums over all of them.
% - Some representations give different molecules identical feature vectors.
% - In the exploratory study of learning with noisy labels, configurations that used a distance between molecules did worst.
% - Judging a label by its neighbours usually made the model worse.
% - What a representation would need: features chosen for the task, few of them, from several representations, traceable to their source.
% Paragraph 4: the adaptive hybrid representation (Chapter~\ref{chap:method}).
% - Asked: how do you select the best set of features for a given prediction task?
% - The method: choose features from several representations for one task, under a budget in features, with every choice made inside each fold.
% - Result needed: where the hybrid beat the best single representation and where it did not, in each dataset's own metric.
% - Result needed: whether datasets preferred different search settings; transfer, antibody and separate-site results in one line each.
% - Result needed: whether importance weighting helped during hybrid construction and model training, in each dataset's own metric.
% - Limits: refer to Chapter~\ref{chap:method}'s limitations section, do not list them.
% Paragraph 5: hybrids for similarity and label noise (Chapter~\ref{chap:similarity}).
% - Asked: does a representation chosen for the task repair nearest-neighbour prediction, noise detection and learning with noisy labels?
% - Result needed: the answer per question, and under which noise types any LNL method helped.
% Paragraph 6: how the pieces fit (short).
% - Chapter~\ref{chap:robustness} tells a reader which model to choose under noise. Chapters~\ref{chap:flaws} to~\ref{chap:similarity} tell them which features to give it.
% - What works depends on the dataset; the next section argues this.
% - Contributions in one line: the NoiseInject package and paper; the adaptive hybrid representation and its blueprints; evidence that distance-based noise handling on standard representations usually made models worse.
% - Then refer to "Summary of Research Contributions" for the papers and software.
% ---- END NOTES ----

\todo[inline, color=yellow!40]{Write Sec 7.1. Draw on [NI 586], already used in the Chapter 3 conclusions -- do not simply repeat it.}


\section{No single method fits every dataset}
% ---- NOTES: No single method fits every dataset ----
% Takeaway: no single setting of the hybrid search, and no single LNL method, was best on every dataset.
% Takeaway: the best choice depends on the noise type and amount, the number of molecules, whether the features describe the property, and the downstream model.
% Paragraph 1: the claim.
% - The search settings are the downstream model, budget, number of base representations, importance method, allocation method and internal evaluator. None was best everywhere.
% - The same holds for label noise: which LNL method helps, if any, depends on the data.
% - LNL methods exist. Noise mitigation, the subset that changes the labels, is the harder part (Chapter~\ref{chap:similarity}).
% - No overall winner averaged across datasets.
% Paragraph 2: what the best choice depends on.
% - The type of label noise (see "By noise type", Chapter~\ref{chap:similarity}).
% - The amount of label noise.
% - The number of molecules. A small dataset leaves few molecules to score each setting on inside a fold, so the chosen setting is less stable (size threshold: see "Limitations", Chapter~\ref{chap:method}).
% - Whether any feature describes the property at all. The search cannot add signal the inputs do not carry.
% - The downstream model. The best hybrid for a tree ensemble need not be the best for a neural network, so the internal evaluator is searched per downstream model.
% Paragraph 3: the evidence, collected (refer to each section rather than restating its numbers).
% - Chapter~\ref{chap:robustness}: the representation with the highest \aucnorm for a model changed between Gaussian and grouped-shifted noise, for five of the thirteen base models on QM9.
% - Result needed: confirm the count of five against the run the thesis reports.
% - Chapter~\ref{chap:robustness}: model ranks moved with the amount of noise. On ECFP4 under Gaussian noise, RF moved from mid-ranking on clean labels to the top at higher noise levels.
% - Chapter~\ref{chap:robustness}: tree and kernel models kept a similar \aucnorm on every representation. The plain and Bayesian networks and NGBoost depended on the representation.
% - Chapter~\ref{chap:robustness}: on the three assay datasets, XGBoost and LightGBM ranked lower for robustness than on QM9.
% - Chapter~\ref{chap:method} (preliminary runs): a fixed share of features per representation won some combinations of dataset, downstream model and budget; one pooled ranking across all representations won others.
% - Chapter~\ref{chap:similarity}, censoring (preliminary): on hERG Ki with 10% of training labels clipped, raising every clipped label by one constant lowered RMSE on all four downstream models tried, on ECFP4.
% - Chapter~\ref{chap:similarity}, Gaussian noise (preliminary): on antibody thermostability, how many of the twelve LNL methods beat the plain network changed with the amount of noise and with the input.
% - Result needed: per dataset, the winning allocation method, number of base representations, importance method and budget.
% - Result needed: per dataset and downstream model, the winning internal evaluator.
% - Result needed: per dataset, whether a fixed share per representation or one pooled ranking won.
% - Result needed: which LNL methods, including the noise-mitigation ones, helped under each noise type.
% Paragraph 4: what the thesis offers instead of one recipe.
% - Not a fixed recipe or a new architecture tuned to one problem.
% - A set of methods that might work, and a procedure that gives one supervised task its best chance by choosing its features.
% - For a practitioner: there is no default to trust. The useful question is which settings to try and how to compare them fairly.
% Figure: a grid, one row per dataset, one column per search setting, each entry coloured by the winning value. Shows whether any column is one colour. Needs the full runs.
% ---- END NOTES ----


\subsection{Simple is often better}
% ---- NOTES: Simple is often better ----
% Takeaway: across the thesis, the simpler choice often did as well as the complex one, or better.
% Takeaway: four places show it: the model under label noise, the number of features, nearest-neighbour prediction, and LNL methods. No new results here.
% Paragraph 1: simple models under label noise (Chapter~\ref{chap:robustness}; robustness is \aucnorm).
% - On QM9 under Gaussian noise, forests, boosted trees, SVM and the GP were both accurate on clean labels and robust.
% - The four plain and Bayesian networks were among the most accurate on clean labels and the four least robust (\aucnorm 0.881 to 0.897, ECFP4).
% - Their \aucnorm also varied more across the six representations than any other model's.
% - Word it as "tree ensembles and SVM, not plain neural networks": the variational networks sat in the middle (\aucnorm 0.924 and 0.926).
% - Robust is not accurate. NGBoost had the highest \aucnorm on all six representations under Gaussian noise on QM9, and the lowest clean R2 on four of the six.
% - Rankings did not carry over whole to the assay datasets. NGBoost ranked 1st on QM9 at every representation, but 6th to 9th on hERG Ki.
% - Random forest ranked 1st to 7th on every assay dataset and representation, and rises as noise grows.
% Paragraph 2: fewer features is often better (Chapters~\ref{chap:method} and~\ref{chap:similarity}).
% - The search sweeps the budget in features and does not assume more features are better.
% - Result needed: how often the best hybrid uses fewer features than the largest budget, per dataset and downstream model.
% - Result needed: whether an ensemble of single-representation models matches the hybrid at lower cost, per dataset.
% - Predicting each molecule's activity from its five nearest neighbours: a small, task-relevant hybrid against each downstream model's top two single representations.
% - Result needed: that comparison per dataset, in the dataset's own metric.
% - The model stays simple; any gain comes from the input.
% Paragraph 3: LNL methods that did not help (no method-level detail).
% - Methods that used a distance between molecules in the exploratory study (Section~\ref{sec:lnl_motivation}), and neighbour-based label cleaning ("Noise detection by neighbours").
% - The twelve LNL methods compared with one plain neural network under added Gaussian noise ("Gaussian noise: learning-with-noisy-labels methods").
% - Result needed: at 0.5 and 1.5 times the labels' standard deviation, how many of the twelve beat the plain network, per dataset, on each model's top two single representations and on the hybrid.
% - Cleaning real labels, nothing injected ("Real labels, nothing injected: cleaning").
% - Result needed: per dataset, whether any cleaning method scored better than leaving the labels as they came.
% - Labels clipped at a detection limit are a separate case where correction helped. Do not write that every method failed.
% - Result needed: whether removing flagged molecules at separate sites with different noise levels improved on pooling the sites.
% Paragraph 4: why LNL methods from other fields do not carry over to assay data.
% - Most LNL methods were built and tested on large image datasets with flipped class labels \citep{Song2022}.
% - Assay datasets are small (1,415 to 5,039 molecules in the three assay datasets of Chapter~\ref{chap:robustness}). Their labels are continuous, and every label carries some error.
% - Many methods rely on neighbours, or on how easily a sample is learned. Both need a space where close molecules have similar activity; Chapter~\ref{chap:flaws} shows standard representations do not give this.
% - On medium-sized tabular data, tree models still beat deep networks, partly because uninformative features hurt networks more \citep{Grinsztajn2022}. One sentence.
% - Financial price data: \citet{Ma2022} treat labels built from price movements as noisy throughout, because one cannot tell whether a price rose from signal or noise.
% - They replace those labels with labels from a denoising autoencoder, and report better downstream classification. Their labels are still classes (price up or down).
% Paragraph 5: tuning a method to one problem.
% - A method tuned closely to one dataset or task loses generality.
% - The tasks here span binding affinity, ADME, quantum properties and antibodies. That range is why learned models are used at all.
% - Simple methods do not generalise automatically either: feature selection needs tuning per dataset, which is why the search sweeps its settings.
% - The search can over-fit too (Discussion of Chapter~\ref{chap:method}).
% - The balance: fix the simple, robust parts (a tree ensemble, or neighbours on a small feature set); search only the choices that vary by dataset.
% ---- END NOTES ----


\section{Limitations}
% ---- NOTES: Limitations ----
% Takeaway: correcting label noise without ground truth is hard, with no clear general solution.
% Takeaway: prefer LNL methods that change training over noise mitigation, except where the noise has a known structure.
% Paragraph 1: correcting label noise without ground truth.
% - Limits of the hybrid search itself (run time, repeating the search, reproducibility) are in "Limitations" in Chapter~\ref{chap:method}. One sentence.
% - The thesis set out to mitigate label noise. The conclusion: it is a hard problem with no clear general solution.
% - Every experimental label carries an unknown amount of noise. There is no clean copy to check a correction against.
% - Finding a noisy label is easier than correcting it (see "By noise type").
% - What is possible: outlier detection, and detection of systematic errors with a known structure (a shared shift per source, clipping at a detection limit).
% - In the assay experiments, "clean" means "no added noise", not "true". Only QM9 is close to noise-free.
% - The injected noise conditions are models of noise. Real data mix them in unknown amounts, so a method that works under one condition may fail on real data.
% - Result needed: how many LNL methods, and how many noise-mitigation ones, improved on the uncorrected model, per noise condition, and where recovery went negative.
% Paragraph 2: recommendation.
% - Noise mitigation estimates each label's error and writes the correction into the data before training. A wrong guess stays in the data for every later model.
% - LNL methods that change training adapt throughout training and use all the data \citep{Song2022}.
% - Noise mitigation is cheaper: done once, reused by any model.
% - Correcting labels directly is justified with a known direction (censoring: the limit and direction are known), or a clearly separate subset (gross outliers, an offset tied to one source).
% - Pair either approach with uncertainty estimation where it is reliable. The overall uncertainty level rises with the amount of noise (Chapter~\ref{chap:robustness}).
% - Per-molecule uncertainty does not pick out which labels are noisy. Use uncertainty to judge how noisy a dataset is, not which labels are wrong.
% Paragraph 3: other limits of scope.
% - Classification: the label-noise experiments are regression only. Flipped class labels are the first next step.
% - Datasets: the controlled label-noise experiments use QM9 and three small assay datasets (LogD, Caco-2, hERG).
% - No dataset has repeated measurements with a known true value for every molecule. Testing label correction directly would need one.
% - Compute limited how many noise conditions, noise levels, folds and datasets could be crossed.
% - Results are per-fold scores and their spread over 3 to 5 folds, with no significance tests. Some patterns are a direction only.
% - Result needed: final number of folds and number of datasets per chapter.
% ---- END NOTES ----
\todo[inline, color=yellow!40]{7.5: noise mitigation is hard; recommend LNL over direct mitigation. Method limitations are in 5.10.}


\section{Future directions}
% ---- NOTES: Future directions ----
% Takeaway: the first next step is classification with flipped labels; then hybrids as inputs to generative models, and a faster search.
% Paragraph 1: classification and label flipping.
% - The noise work here is almost all regression. In classification, a noisy label is a flipped class, not a shifted value.
% - The exception: the survey of twelve LNL methods flips labels uniformly at random on its classification datasets. Ten of the twelve methods run there.
% - Many QSAR class labels come from cutting a continuous value at a threshold. Flips then concentrate near the cut-off, so the noise depends on the molecule.
% - NoiseInject already has six flip conditions: uniform, class imbalance, binary asymmetric, instance-dependent, class-dependent and confusion-directed. The hybrid search already handles classification.
% - Missing: a study of the molecule-dependent flip conditions at the scale of the regression study.
% - Open question: do the regression findings carry over? Most LNL methods that change training were designed for flipped labels.
% - Limit: nearest-neighbour prediction is scored by AUROC, defined for two classes. It gives no score with more than two classes, or on a held-out block that holds one class.
% Paragraph 2 (her paragraph "I also plan to extend..."): other representations and importance methods.
% - Change: cut GNN embeddings and sequence encoders; the search already includes them.
% - Change: cut TreeSHAP and gradient-based attribution; TreeSHAP and DeepLIFT are already swept, with LightGBM gain.
% - Change: cut small-loss selection; it scores training molecules, not features, and is one of the twelve LNL methods surveyed.
% - Still open: transformed inputs (normalised or binarised continuous features) and features derived from graph distances. Neither is searched.
% - Still open: importance methods that treat the model as a black box, e.g. KernelSHAP, or dropping one feature and refitting.
% - The software already has KernelSHAP and permutation importance, but the search sweeps neither. Dropping a feature and refitting is not implemented.
% - Change: keep "and datasets" only if a domain outside those in "Datasets and splitting" is named.
% Paragraph 3 (her paragraph "As mentioned earlier..."): model against representation complexity.
% - Change: cut the plan to adapt the hybrid to datasets and models. The search already keeps one best configuration per dataset, downstream model and budget.
% - Change: cut the planned trade-off study. Chapter~\ref{chap:method} reports held-out score against budget per model, and tests whether the internal evaluator must match the downstream model.
% - Change: cut "a very small default hybrid with RF-based importance"; it describes only the noise detection study (a hybrid of 50-100 features), not the current search.
% - Change: cut "I do not currently have a full set of results on hybrid"; Chapters~\ref{chap:method} and~\ref{chap:similarity} carry them.
% - Keep one sentence: a search objective that rewards robustness to label noise as well as accuracy. It is not built.
% Paragraph 4 (her paragraph "Various estimation methods..."): other noise types.
% - Chapter~\ref{chap:similarity} already compares LNL methods under censoring, Gaussian noise and real labels.
% - Still open: LNL methods under outliers, and under grouped offsets (a whole laboratory or source document shifted together). Chapter~\ref{chap:similarity} uses both only to grade noise detection.
% - Keep: whether cleaning based on each molecule's nearest neighbours, which fails under Gaussian noise, does better on outliers.
% - Keep: whether label smoothing and sample weighting still lead under those noise types.
% Paragraph 5: hybrid features as inputs to generative models.
% - The model that scores molecules for a molecule generator could take the hybrid as its input.
% - Open question: what should the search optimise for that task? Held-out accuracy on known molecules may not reward the right thing for new ones.
% - Cost: every generated molecule needs its hybrid features computed, including a pass through any embedding model in the hybrid.
% Paragraph 6: a faster search.
% - As it stands, the full search is too heavy for routine use.
% - Two ways to cut cost exist: the fast tier searches shorter value lists (Table~\ref{tab:search_tiers}), and adaptive pruning drops configurations early.
% - Next step: a warm start. Begin from the blueprint of a similar dataset and search near it.
% - Chapter~\ref{chap:method} already reuses one dataset's search settings on a second dataset with no search. A warm start would search near those settings instead.
% - A faster search must still be able to reach every configuration. It only visits fewer of them.
% - Result needed: hours and number of model fits for the full search against the fast tier, per dataset.
% Paragraph 7: one or two sentences each.
% - Forward compatibility both ways: does the hybrid improve a weak new embedding, not only absorb a strong one?
% - An experiment scores the hybrid that includes a new embedding against the embedding alone with its best downstream model.
% - It also scores it against a hybrid where a random matrix with the same number of features replaces the embedding.
% - Blueprints across domains, at the level of representation families. Individual features cannot transfer between domains, because the representations differ.
% - Data held at separate sites: Chapter~\ref{chap:method} used the hybrid indirectly with such data.
% - Next step: a fuller federated-learning deployment with Lhasa, with real sites running the selection on their own data \citep{McMahan2017}.
% Paragraph 8 (the paragraph "What does this mean in practice?"): adding a new foundation model.
% - Change: say "the hybrid search" instead of "KIRBy", which names the repository only; say "the representations searched" instead of "the pool".
% - Change: "a substantive change" and "a principled combination" claim the method's value; tone them down or tie them to Chapter~\ref{chap:method}'s results.
% Figure: hub-and-spoke, with the hybrid search in the centre. Each spoke is tagged with the chapter it extends.
% - Spokes: classification, generative scoring, sites holding their own data, a faster search, and new embeddings or domains.
% ---- END NOTES ----

I also plan to extend the hybrid research to additional base representations and datasets. Some alternative representations include additional GNN embeddings, sequential encoders, augmented features such as embedded distances from molecular graphs, and transformed representations (i.e., normalization, binarization of continuous features). Additionally, I want to explore alternative feature importance scoring methods, especially ones that involve different families of models. Specifically, I plan to look at TreeSHAP, KernelSHAP, drop-column importance for noise robustness, gradient-based attribution for NN embeddings, and small-loss selection, as discussed in Section~\ref{sec:lnl_motivation}.

As mentioned earlier, although I do not currently have a full set of results on hybrid, they have been showing promise throughout a rank of QSAR and noise mitigation tasks. I would like to explore the complexity tradeoff between model and representation complexity, as I've seen a trend of complicated models like NNs having stronger predictive performance with smaller hybrids and more simplistic tree-based methods benefiting from larger hybrids. This noise mitigation experiment relied on a very small default hybrid and used RF-based feature importance. I would like to leverage my allocation strategies to adapt hybrid models to datasets \textit{and} models. These strategies can also be modified to optimize for noise robustness in addition to predictive performance.

Various estimation methods showed promise, meriting further investigation. Additionally, neighbor-based noise mitigation strategies which are not fit for Gaussian noise may perform differently on outlier-based strategies. The dominance of estimators with label smoothing and sample weighting may be challenged for different types of noise.

What does this mean in practice? A new foundation model appears every few months in molecular, protein, or antibody modeling. A researcher using KIRBy does not need to decide whether to replace their current encoding or retrain their pipeline; they compute the new model's embeddings, add the matrix to the pool, and re-run the search. The search determines whether the new embedding contributes useful features and, if so, how many and which ones. In a field where encoding choice is often determined by familiarity or lab convention rather than by systematic comparison, this is a substantive change. It replaces a single bet with a principled combination, tailored to each task, that updates as the pool grows.


\pagebreak
\chapter*{Summary of Research Contributions}
\addcontentsline{toc}{chapter}{Summary of Research Contributions}

% ---- NOTES: Summary of Research Contributions ----
% Takeaway: a short list of the papers, the conference presentation, the blog posts and the software written for this thesis.
% Paragraph 1: order.
% - Research papers, one entry per paper, in chapter order; then the conference presentation; then the blog posts.
% - Then a new \section{Software}, after the blog posts. Every entry is a citation-style line and at most one or two sentences.
% Software, opening sentence: three pieces of open-source software were written for this thesis.
% - Each entry: the name, one sentence on what it computes, and its link. Licence, version and requirements: refer once to Appendix~\ref{app:availability}.
% Software entry 1: NoiseInject (released).
% - A Python package that adds label noise to training labels and scores how a model's performance and uncertainty change as the noise grows.
% - It builds the noise conditions of Chapter~\ref{chap:robustness} and computes the robustness and uncertainty metrics defined there.
% - Link github.com/adpunt/noise_inject; archived on Zenodo at https://doi.org/10.5281/zenodo.20532710; cite \citep{noiseinject}.
% Software entry 2: the QM9 pipeline (the Rust part to be released).
% - Public repository github.com/adpunt/qsar_qm_models, GPL-3.0 licence.
% - A Python program computes the six representations and the scaffold groups.
% - A Rust program adds label noise to the training and validation labels, records the noise drawn for each molecule, and writes memory-mapped files for the models.
% - It produced the QM9 results of Chapter~\ref{chap:robustness}.
% Software entry 3: KIRBy (to be released).
% - The Python research repository: the adaptive hybrid representation search, the representation generators, the dataset loaders, and the experiments of Chapters~\ref{chap:flaws} to~\ref{chap:similarity} and the assay-dataset runs of Chapter~\ref{chap:robustness}.
% - It contains no Rust. The hybrid search works on feature matrices, so it runs without the QM9 pipeline.
% ---- END NOTES ----

\section{Research Papers}
% ---- NOTES: Research Papers ----
% Takeaway: one entry per paper, in chapter order: the NoiseInject paper, then two hybrid papers.
% Paragraph 1: the NoiseInject paper (Chapter~\ref{chap:robustness}).
% - Change: title to "NoiseInject: A Benchmarking Framework to Study Noise Robustness in QSAR Models".
% - Change: author order to Punt, Hanser, Werner, Morris.
% - Target stays the Journal of Cheminformatics; status stays "in preparation".
% - Change: "prepration" to "preparation", in both entries.
% Paragraph 2: the hybrid work, as two papers.
% - Change: split the second entry into two, each with its own title and one sentence on its status.
% - Paper on building the adaptive hybrid representation (Chapter~\ref{chap:method}): when and why the hybrid helps, and whether it is explainable.
% - It also asks whether the hybrid extends to new embeddings and to data held at separate sites.
% - Paper on hybrid representations for molecular similarity and label noise (Chapter~\ref{chap:similarity}): whether the hybrid gives a more property-aware similarity and whether it improves noise estimation.
% - Change: the current title says "noise prediction robustness and mitigation"; Chapter~\ref{chap:similarity} covers noise detection and LNL, of which noise mitigation is one part.
% - Working titles and venue for both: not yet chosen.
% ---- END NOTES ----
\textbf{In prepration:} Punt, A., Morris, G., Werner, S., \& Hanser, T. (2026). Noise-Robust QSAR: evaluating probabilistic methods across molecular representations \\
My supervisors and I are currently iterating on an advanced draft for this research paper based on the material in Chapter~\ref{chap:robustness}, which we intend to submit to the \textit{Journal of Cheminformatics} within the next few months.

\textbf{In prepration:} Punt, A., Morris, G., Werner, S., \& Hanser, T. (2026). Improving noise prediction robustness and mitigation with feature-importance hybrid representations in multi-modal contexts \\ I am in the early stages of writing a draft and extending the experiments in Chapters~\ref{chap:method} and~\ref{chap:similarity} within the next few months.

\section{Conference Presentations}
% ---- NOTES: Conference Presentations ----
% Takeaway: the one conference entry, in the same citation layout as the papers.
% - Change: "Posted presented" to "Poster presented".
% ---- END NOTES ----
Punt, A., Morris, G., Werner, S., \& Hanser, T. (2024). Investigating the impact and mitigation of noise in QSAR models. Posted presented at the Autumn UK-QSAR 2024 Meeting.


\section{Blogposts}

I have contributed to the \textit{Blopig} blog of the Oxford Protein Informatics Group (OPIG):

\begin{itemize}
    \item \href{https://www.blopig.com/blog/2025/10/are-you-addicted-to-dopamine/}{Are you addicted to dopamine?}
    \item \href{https://www.blopig.com/blog/2025/05/demystifying-git-and-merge-conflicts/}{Demystifying Git and merge conflicts}
    \item \href{https://www.blopig.com/blog/2025/03/ai-in-academic-writing-ethical-considerations-and-best-practices/}{AI in Academic Writing: Ethical Considerations and Best Practices}
    \item \href{https://www.blopig.com/blog/2024/08/memory-mapped-files-for-efficient-data-processing/}{Memory-mapped files for efficient data processing}
    \item \href{https://www.blopig.com/blog/2024/04/exploring-multilingual-programming/}{Exploring multilingual programming}
\end{itemize}

%%========================= APPENDICES ======================================

\appendix

